\section{Probabilidade e Estatística}

\questao{(OM$^3$SP 2026 - 1$^\circ$ fase ) Os professores de Matemática da Escola Marcelo Minervo dos Santos Pimentel, realizaram uma pesquisa com 480 pessoas. Os professores perguntava aos entrevistados qual gênero de filme eles preferiam: ação, comédia, anime, terror ou outros. Este gráfico mostra o percentual de pessoas que responderam à pesquisa.}

        \begin{tikzpicture}[scale=.7]
            \pie[
                text=legend,
                color={red!60, yellow!75, cyan!80, orange!75, blue!50}
            ]{5/Ação, 35/Comédia, 20/Anime, 25/Terror, 15/Outros}
        \end{tikzpicture}

        Quantas pessoas responderam que preferem o gênero de filme anime?


            \begin{enumerate}[label=\alph*)]
                \item 96 pessoas
                \item 105 pessoas
                \item 90 pessoas
                \item 85 pessoas
            \end{enumerate}

\begin{solucoes}{1}
    \solucao
    Para determinar a quantidade de pessoas que preferem o gênero de filme anima, observaremos qual a sua porcentagem no gráfico e depois calcularemos quanto equivale do total de pessoas entrevistadas, assim, temos:

$\begin{WithArrows}[tikz={>=stealth, text width = 6cm, draw=black!70, text=black!70}, interline=2mm]
    Anime =& 20\%\, de\, 480  \Arrow[rl]{Transformando em um produto} \\ 
    Anime =& \dfrac{20}{100} \times 480   \\
    Anime =& \dfrac{\cancel{20}}{\cancel{100}} \times 480 
    \Arrow[ll]{Simplificando a fração por 20, temos}\\
    Anime =& \dfrac{480}{4} \\
    Anime =& 96
\end{WithArrows}$

Sendo assim, a quantidade de pessoas que preferem o gênero de filme Anime é 96

 \textbf{ALTERNATIVA D}

\end{solucoes}

\questao{Para escolher 1 entre 4 estudantes para ser presidente do grêmio da Escola Marcelo Pimentel, o coordenador Jefferson estabeleceu como critério a maior média obtida com as notas dadas \textbf{à prova escrita}, \textbf{à entrevista} e a \textbf{apresentação de ideias}. Analise as notas obtidas pelos estudantes e descubra qual o estudante escolhido.}

        \begin{figure}[H]
        \centering
            \begin{minipage}{.33\linewidth}
                \begin{tabular}{|p{2.5cm}|}
                    \hline
                    \multicolumn{1}{|c|}{ESTUDANTE 1}\\
                    \hline
                    Prova: 8 \par
                    Entrevista: 7 \par 
                    Ideias: 6
                    \\
                    \hline
                    
                \end{tabular}
            \end{minipage}
            \begin{minipage}{.33\linewidth}
                \begin{tabular}{|p{2.5cm}|}
                    \hline
                    \multicolumn{1}{|c|}{ESTUDANTE 2}\\
                    \hline
                    Prova: 8 \par
                    Entrevista: 8 \par 
                    Ideias: 7
                    \\
                    \hline
                    
                \end{tabular}
            \end{minipage}
            
        \end{figure}

\vspace*{-15pt}
        \begin{figure}[H]
            \centering
            \begin{minipage}{.33\linewidth}
                \begin{tabular}{|p{2.5cm}|}
                    \hline
                    \multicolumn{1}{|c|}{ESTUDANTE 3}\\
                    \hline
                    Prova: 7 \par
                    Entrevista: 6 \par 
                    Ideias: 7
                    \\
                    \hline
                    
                \end{tabular}
            \end{minipage}
            \begin{minipage}{.33\linewidth}
                \begin{tabular}{|p{2.5cm}|}
                    \hline
                    \multicolumn{1}{|c|}{ESTUDANTE 4}\\
                    \hline
                    Prova: 8 \par
                    Entrevista: 9 \par 
                    Ideias: 5
                    \\
                    \hline
                    
                \end{tabular}
            \end{minipage}
           
        \end{figure}

        \begin{multicols}{2}
            \begin{enumerate}[label=\alph*)]
                \item ESTUDANTE 1
                \item ESTUDANTE 2
                \item ESTUDANTE 3
                \item ESTUDANTE 4
            \end{enumerate}
        \end{multicols}

\begin{solucoes}{1}
\solucao

Para descobrir qual estudante foi escolhido, é preciso calcular a média de cada um dos estudantes.

\begin{multicols}{2}
    Para o estudante 1, temos:

$\begin{WithArrows}[tikz={>=stealth, text width= 6cm, draw=black!70, text=black!70}, interline=2mm]
   M &= \dfrac{8+7+6}{3} \\
   M &= \dfrac{21}{3} \\
   M &= 7
\end{WithArrows}$

Média igual a 7.

Para o estudante 2, temos:

$\begin{WithArrows}[tikz={>=stealth, text width = 6cm, draw=black!70, text=black!70}, interline=2mm]
   M &= \dfrac{8+8+7}{3}\\
   M &= \dfrac{23}{3} \Arrow[ll]{Calculando o resultado da fração, temos aproximadamente}\\
   M &= 7,6
\end{WithArrows}$

Média igual a 7,6.

Para o estudante 3, temos:

$\begin{WithArrows}[tikz={>=stealth, text width = 6cm, draw=black!70, text=black!70}, interline=2mm]
    M &= \dfrac{7+6+7}{3} \\
    M &= \dfrac{20}{3} \Arrow[ll]{Calculando o resultado da fração, temos aproximadamente}\\
    M &= 6,6
\end{WithArrows}$

Média igual a 6,6.

Para o estudante 4, temos:

$\begin{WithArrows}[tikz={>=stealth, text width = 6cm, draw=black!70, text=black!70}, interline=2mm]
    M &= \dfrac{8+9+5}{3} \\
    M &= \dfrac{22}{3} \Arrow[ll]{Calculando o resultado da fração, temos aproximadamente}\\
    M &= 7,3
\end{WithArrows}$

Média igual a 7,3.
\end{multicols}

Fazendo a comparação das médias obtidas por cada estudante, temos:

estudante 2 $>$ estudante 4 $>$ estudante 1 $>$ estudante 3.

Assim, o escolhido foi o \textbf{estudante 2}

\textbf{ALTERNATIVA B}
    
\end{solucoes}